%------------------------------------------------------------------------------%
\documentclass[a4paper,12pt,fleqn]{article}

\usepackage{calculo_varias_variaveis-1}

\ShowAnswers

%------------------------------------------------------------------------------%
\begin{document}

\makeHeader{CFVVI}{Prova 1 -- Turma 2}{22/09/2025}

% 01
%------------------------------------------------------------------------------%
\begin{multicols}{2}
\raggedbottom
\exercise{50\%}
  Uma partícula descreve uma trajetória elíptica descrita
  pelas equações paramétricas
  \begin{gather*}
    x(t) = 5 \sin\left(\frac{\pi}{2}-t\right) - 1\\
    y(t) = 3 \cos\left(\frac{\pi}{2}-t\right) + 1\\
    0 \leq t \leq \pi
  \end{gather*}
  \begin{enumerate}[label=\alph*),itemsep=0pt,topsep=0pt]
    \item\!\! [20\%] calcule o vetor tangente à curva
    \item\!\! [10\%] avalie o vetor tangente em \\
           \hspace*{29pt} $t=0$, $t=\frac{\pi}{2}$ e em $t=\pi$
    \item\!\! [20\%] esboce a curva e os vetores \\
           \hspace*{29pt} tangentes calculados
  \end{enumerate}
\columnbreak
~

\begin{questiononly}
  \hfill\includegraphics{axis_xy}
\end{questiononly}
\begin{answeronly}
  \hfill\includegraphics{T2_B_parametrica}
\end{answeronly}
\end{multicols}
\clearpagequestiononly

\begin{answer}
\textbf{a)} \quad
O vetor tangente é
\(
  v(t) = \Vetor{\frac{dx}{dt}}{\frac{dy}{dt}}
\)

Avaliando cada derivada
\medskip
\[
  \frac{dx}{dt}
  = \frac{d}{dt}\left[5 \sin\left(\frac{\pi}{2}-t\right) - 1\right]
  = 5 \frac{d}{dt}\sin\left(\frac{\pi}{2}-t\right)
  = 5 \cos\left(\frac{\pi}{2}-t\right)\frac{d}{dt}\left(\frac{\pi}{2}-t\right)
  = -5 \cos\left(\frac{\pi}{2}-t\right)
\]

\medskip
\[
  \frac{dy}{dt}
  = \frac{d}{dt}\left[3 \cos\left(\frac{\pi}{2}-t\right) + 1\right]
  = 3 \frac{d}{dt}\cos\left(\frac{\pi}{2}-t\right)
  = -3 \sin\left(\frac{\pi}{2}-t\right)\frac{d}{dt}\left(\frac{\pi}{2}-t\right)
  = 3 \sin\left(\frac{\pi}{2}-t\right)
\]

\medskip
Portanto,
\(
  v(t) = \Vetor{-5 \cos\left(\frac{\pi}{2}-t\right)}{3 \sin\left(\frac{\pi}{2}-t\right)}
\)

\clearpage
\textbf{b)} \quad
Avaliando o vetor tangente nos pontos $t=0$, $t=\frac{\pi}{2}$ e em $t=\pi$
\setlength{\mathindent}{0pt}
\begin{multicols}{3}
  \begin{align*}
    v(0)
    & = \Vetor{-5 \cos\left(\frac{\pi}{2}-t\right)}{3 \sin\left(\frac{\pi}{2}-t\right)} \evalat{0}{} \\[2mm]
    & = \Vetor{-5 \cos\left(\frac{\pi}{2}\right)}{3 \sin\left(\frac{\pi}{2}\right)} \\[2mm]
    & = \vetor{-5 \times 0}{3 \times 1} \\[2mm]
    & = \vetor{0}{3}
  \end{align*}

  \newcolumn
  \begin{align*}
    v\left(\frac{\pi}{2}\right)
    & = \Vetor{-5 \cos\left(\frac{\pi}{2}-t\right)}{3 \sin\left(\frac{\pi}{2}-t\right)} \evalat{\frac{\pi}{2}}{} \\[2mm]
    & = \Vetor{-5 \cos\left(0\right)}{3 \sin\left(0\right)} \\[2mm]
    & = \Vetor{-5 \times 1}{3\times 0} \\[2mm]
    & = \vetor{-5}{0}
  \end{align*}

  \newcolumn
  \begin{align*}
    v(\pi)
    & = \Vetor{-5 \cos\left(\frac{\pi}{2}-t\right)}{3 \sin\left(\frac{\pi}{2}-t\right)} \evalat{\pi}{} \\[2mm]
    & = \Vetor{-5 \cos\left(\frac{\pi}{2}-\pi\right)}{3 \sin\left(\frac{\pi}{2}-\pi\right)} \\[2mm]
    & = \Vetor{-5 \cos\left(-\frac{\pi}{2}\right)}{3 \sin\left(-\frac{\pi}{2}\right)} \\[2mm]
    & = \vetor{-5 \times 0}{3(-1)} \\[2mm]
    & = \vetor{0}{-3}
  \end{align*}
\end{multicols}

\bigskip
\textbf{b)} \quad
Calculando os pontos associados a $t=0$, $t=\frac{\pi}{2}$ e em $t=\pi$
\[
  x(0)
  = \left(5 \sin\left(\frac{\pi}{2}-t\right) - 1\right)\evalat{0}{}
  = 5 \sin\left(\frac{\pi}{2}\right) - 1
  = 5 \times 1 - 1
  = 4
\]
\[
  y(0)
  = \left(3 \cos\left(\frac{\pi}{2}-t\right) + 1\right)\evalat{0}{}
  = 3 \cos\left(\frac{\pi}{2}\right) + 1
  = 3 \times 0 + 1
  = 1
\]

\medskip
\[
  x\left(\frac{\pi}{2}\right)
  = \left(5 \sin\left(\frac{\pi}{2}-t\right) - 1\right)\evalat{\frac{\pi}{2}}{}
  = 5 \sin\left(0\right) - 1
  = 5 \times 0 - 1
  = -1
\]
\[
  y\left(\frac{\pi}{2}\right)
  = \left(3 \cos\left(\frac{\pi}{2}-t\right) + 1\right)\evalat{\frac{\pi}{2}}{}
  = 3 \cos\left(0\right) + 1
  = 3 \times 1 + 1
  = 4
\]

\medskip
\[
  x(\pi)
  = \left(5 \sin\left(\frac{\pi}{2}-t\right) - 1\right)\evalat{\pi}{}
  = 5 \sin\left(-\frac{\pi}{2}\right) - 1
  = 5 (-1) - 1
  = -6
\]
\[
  y(\pi)
  = \left(3 \cos\left(\frac{\pi}{2}-t\right) + 1\right)\evalat{\pi}{}
  = 3 \cos\left(-\frac{\pi}{2}\right) + 1
  = 3 \times 0 + 1
  = 1
\]
\end{answer}

% 02
%------------------------------------------------------------------------------%
\exercise{50\%}
Considerando a função
\(
  f(x, y) = \frac{2}{\sqrt{x^2 + y^2 -4}}
\)
\begin{enumerate}[label=(\alph*),itemsep=0pt]
  \item \ [10\%] determine o domínio de $f$
  \item \ [10\%] esboce o domínio
  \item \ [20\%] determine todas as curvas de nível de $f$
  \item \ [10\%] determine e esboce a curva de nível que passa pelo ponto \((2, 4)\)
\end{enumerate}
\begin{questiononly}
\begin{multicols}{2}
  \;\; Domínio
  \vspace{-0.3\baselineskip}
  \begin{center}
    \includegraphics{axis_xy}
  \end{center}
  \;\; Curva de nível
  \vspace{-0.3\baselineskip}
  \begin{center}
    \includegraphics{axis_xy}
  \end{center}
\end{multicols}
\end{questiononly}
\begin{answeronly}
  \begin{multicols}{2}
    \;\; Domínio
    \vspace{-0.3\baselineskip}
    \begin{center}
      \includegraphics{T2_B_dominio}
    \end{center}
    \;\; Curva de nível
    \vspace{-0.3\baselineskip}
    \begin{center}
      \includegraphics{T2_B_curva_nivel}
    \end{center}
  \end{multicols}
\end{answeronly}

\begin{answer}
\textbf{a)} \quad
  Não podemos calcular a raiz quadrada de um número negativo e não
  podemos dividir por zero, portanto, o domínio de $f$ é o conjunto
  de todos os pontos \((x, y)\in\R^2\) tais que
  \[
    x^2 + y^2 - 4 > 0
  \]
  Podemos reescrever essa condição como
  \[
    x^2 + y^2 > 2^2
  \]
  ou seja, os pontos exteriores ao disco centrado em \((0, 0)\) de raio 2.
  Assim, o domínio é
  \[
    D_f = \{(x, y) \in \R^2 \st x^2 + y^2 > 2^2\}
  \]

\bigskip
\textbf{c)} \quad
  Para determinarmos todas as curvas de nível de $f$ precisamos determinar sua
  imagem. Observamos que dentro do domínio temos
  \begin{align*}
    x^2 + y^2 - 4                  & > 0 \\[2mm]
    \sqrt{x^2 + y^2 - 4}           & > 0 \\[2mm]
    \frac{\sqrt{x^2 + y^2 - 4}}{2} & > 0 \\[2mm]
    0 < \frac{2}{\sqrt{x^2 + y^2 - 4}} & < \infty
  \end{align*}
  Assim, \(\Im_f = \left\{z\in\R \st z > 0\right\}\).
  As curvas de nível de $f$ consistem dos pontos que satisfazem
  \(
    f(x, y) = c
  \)
  para cada \(c\in\Im_f\).
  \begin{align*}
    f(x, y)                        & = c                 \\[2mm]
    \frac{2}{\sqrt{x^2 + y^2 - 4}} & = c                 \\[2mm]
    \sqrt{x^2 + y^2 - 4}           & = \frac{2}{c}       \\[2mm]
    x^2 + y^2 - 4                  & = \frac{4}{c^2}     \\[2mm]
    x^2 + y^2                      & = \frac{4}{c^2} + 4
  \end{align*}
  As curvas de nível são circunferências centradas na origem com raio
  \(
    r
    = \sqrt{\frac{4}{c^2}+4}
    = \frac{2}{c}\sqrt{1+c^2}
  \)

\bigskip
\textbf{d)} \quad
A curva de nível que passa pelo ponto $(2, 4)$ tem o valor
\[
  c
  = f(2, 4)
  = \frac{2}{\sqrt{x^2 + y^2 - 4}}\evalat{(2, 4)}{}
  = \frac{2}{\sqrt{2^2 + 4^2 - 4}}
  = \frac{2}{\sqrt{4 + 16 - 4}}
  = \frac{2}{\sqrt{16}}
  = \frac{2}{4}
  = \frac{1}{2}
\]
Correspondendo a uma circunferência de raio
\[
  r
  = \sqrt{\frac{4}{c^2}+4}\;\evalat{\frac{1}{2}}{}
  = \sqrt{\frac{4}{\sfrac{1}{4}}+4}
  = \sqrt{16+4}
  = \sqrt{20}
  = \sqrt{4\times 5}
  = 2\sqrt{5}
  \approx 4,7321
\]

\end{answer}

%------------------------------------------------------------------------------%
\end{document}
%------------------------------------------------------------------------------%
